\begin{claim}{D.62}{Proved}{Primitive homology centralizers and the exact cusp-transporter root test}
Let $F=\widehat F(x,y)$ and let $g\in F$ have primitive
abelianization: there is a continuous homomorphism
$\rho:F\to\widehat\Z$ with $\rho(g)=1$.
Then, for every nonzero integer $n$,
\[
 C_F(g^n)=\overline{\langle g\rangle}\cong\widehat\Z,
 \qquad
 N_F(\overline{\langle g\rangle})
                  =\overline{\langle g\rangle}.
                                                        \tag{D.62a}
\]
This conclusion uses homology primitivity; it does not
require a free basis containing $g$.

Now use the exact negative torus extension of
D.58--D.61:
\[
 \begin{gathered}
 E=F\rtimes\langle v\rangle,\qquad v^2=1,\qquad
 \alpha=\operatorname{Inn}(v)|_F,\qquad
 c=[x,y],\qquad \alpha(c)=c^{-1},\\
 C=C_F(v),\qquad gvg^{-1}=cv.
 \end{gathered}                                         \tag{D.62b}
\]
D.59 supplies the required primitive abelianization of
this $g$ and the exact inclusion $C\subseteq\overline{F'}$.
For every $g'\in gC$, the following two conditions
are equivalent:
\[
 \begin{split}
 &\alpha(g')\text{ is conjugate to }g'\text{ in }F;\\
 &C_E(g')\text{ contains an element outside }F.
 \end{split}                                           \tag{D.62c}
\]
Their consequences and converse maps are exact:

\begin{enumerate}
\item
In the two-section-class case of D.58,
they are equivalent to the existence of an involution
$w\in E\setminus F$ centralizing $g'$.
That involution is unique, belongs to the other
section class, and
\[
 C_E(g')=\overline{\langle g'\rangle}\times\langle w\rangle.
                                                        \tag{D.62d}
\]
The maps
\[
 X\longmapsto g',\qquad Y\longmapsto vw,\qquad
 v_0\longmapsto v
                                                        \tag{D.62e}
\]
give a full isomorphism from the two-class standard
extension in D.61 to $E$, preserving the quotient labels
and carrying $[X,Y]$ to $c$.

\item
In the one-section-class case of D.58,
they are equivalent to the existence of $z\in E\setminus F$
with $z^2=g'$.
This root is unique, and
\[
 C_E(g')=\overline{\langle z\rangle}\cong\widehat\Z,
 \qquad
 F\cap C_E(g')=\overline{\langle z^2\rangle}
                         =\overline{\langle g'\rangle}.
                                                        \tag{D.62f}
\]
The maps
\[
 X\longmapsto zv,\qquad Y\longmapsto vz,\qquad
 v_0\longmapsto v
                                                        \tag{D.62g}
\]
give the corresponding full isomorphism from the
one-class standard extension, with the same exact
peripheral and quotient labels.
\end{enumerate}

Thus the multiplier-one, $v_0\mapsto v$ portion of the
D.61 conjugator problem is exactly the existence of
a member $g'\in gC$ satisfying (D.62c).
The explicit continuation (D.62q)--(D.62w) treats every
full unit multiplier and both actual reflection labels,
and gives the same criterion for the entire D.61
conjugator space. Neither equivalence asserts that
its required conjugacy occurs.
\end{claim}

\begin{proof}
\emph{The primitive-homology centralizer.}
The equation $\rho(g)=1$ first proves that
$a\mapsto g^a$ embeds $\widehat\Z$ as the indicated
closed procyclic subgroup. Put $D=C_F(g^n)$.
We show $\ker(\rho|_D)=1$.

Suppose $k\ne1$ belongs to this kernel.
The procyclic group generated by $k$ has a nontrivial
$p$-primary component for some prime $p$.
Since $F$ is torsion-free, that component is infinite
procyclic, isomorphic to $\Z_p$.
Let $k_p$ be a topological generator of it.
The $p$-primary component of $g^n$ also generates
a copy of $\Z_p$: under $\rho$ it is $n\Z_p$
in the $p$ coordinate, and $n\ne0$ in $\Z_p$.
These two procyclic groups commute, because $k$
centralizes $g^n$. Their intersection is trivial:
$\rho$ is zero on the first and injective on the
second. Multiplication consequently gives a continuous
injection
\[
                 \Z_p\times\Z_p\longrightarrow F.
                                                        \tag{D.62h}
\]
Its image is closed by compactness. This contradicts
$\operatorname{cd}_p(F)=1$, since a closed subgroup
has no larger $p$-cohomological dimension and
$\operatorname{cd}_p(\Z_p^2)=2$.
Thus $\rho|_D$ is injective.

Every $h\in D$ satisfies
$h g^{-\rho(h)}\in\ker(\rho|_D)$, because $g\in D$.
Hence $h=g^{\rho(h)}$, proving the centralizer equality.
For a normalizing element $h$, write
$hgh^{-1}=g^\lambda$ with
$\lambda\in\widehat\Z^\times$.
Applying $\rho$ gives $1=\lambda$.
The normalizer equality now follows from the
centralizer equality with $n=1$.
No change from homology primitivity to free-basis
primitivity has been made.

\emph{The two exact labels of a central extension.}
Fix $g'\in gC$. Then
\[
 [g',v]=c,\qquad
 \alpha(g')=c^{-1}g',\qquad
 [g']_{\rm ab}=[g]_{\rm ab}.
                                                        \tag{D.62i}
\]
The first equality uses the defining commutation of
$C$ with $v$, and the last uses $C\subseteq\overline{F'}$.
In particular (D.62a) applies to $g'$.
An outside element has the unique form $hv$, $h\in F$,
and
\[
 (hv)g'(hv)^{-1}=h\alpha(g')h^{-1}.
                                                        \tag{D.62j}
\]
This is the actual comparison proving (D.62c), in
both directions.

Suppose its conditions hold and choose
$t\in C_E(g')\setminus F$.
The kernel of the parity map on $C_E(g')$ is exactly
$\overline{\langle g'\rangle}$ by (D.62a).
Thus $t^2=(g')^\kappa$ for a unique
$\kappa\in\widehat\Z$, and $t$ commutes with $g'$.
The entire centralizer has the presentation and
coordinate maps
\[
 \begin{split}
 P_\kappa&=
 \langle G,T\mid [G,T]=1,\ T^2=G^\kappa\rangle_{\rm prof}
          \xrightarrow{\ \sim\ } C_E(g'),\\
 G&\longmapsto g',\qquad T\longmapsto t,\\
 (a,\epsilon)&\longmapsto (g')^a t^\epsilon
       \quad(a\in\widehat\Z,\ \epsilon\in\{0,1\}).
 \end{split}                                           \tag{D.62k}
\]
Surjectivity is the two-coset description.
For injectivity, a relation with $\epsilon=1$
has nonzero parity, whereas a relation with
$\epsilon=0$ is detected by the primitive character
of $g'$. The coordinate multiplication is
\[
 (a,\epsilon)(b,\delta)=
 \bigl(a+b+\epsilon\delta\kappa,\,
                       \epsilon+\delta\bmod2\bigr).
                                                        \tag{D.62l}
\]
These maps also give the specified compact topology.

An outside element $(g')^a t$ is an involution
precisely when $2a+\kappa=0$ in $\widehat\Z$.
This equation has a solution exactly when
$\kappa_2\in2\Z_2$; at every odd prime division
by two is unique. In that case
\[
       w=(g')^{-\kappa/2}t,\qquad
       (a,\epsilon)\longmapsto(g')^a w^\epsilon
\]
is the exact isomorphism
$\widehat\Z\times C_2\to C_E(g')$.
Its outside involution is unique.

If $\kappa_2$ is odd, choose $b\in\widehat\Z$ with
\[
 b_2=0,\qquad
 b_p=(1-\kappa_p)/2\quad(p\text{ odd}),
 \qquad \kappa'=\kappa+2b.
                                                        \tag{D.62m}
\]
The element $\kappa'$ is a full unit: its dyadic
coordinate is $\kappa_2$ and every odd coordinate
is one. Set
\[
 t'=(g')^b t,\qquad
 \nu=(\kappa')^{-1},\qquad z=(t')^\nu.
                                                        \tag{D.62n}
\]
Then $(t')^2=(g')^{\kappa'}$, $z^2=g'$, and $z$
has outside parity because $\nu$ is odd at two.
Moreover $t'=z^{\kappa'}$, so $z$ generates the
entire centralizer. If $z^a=1$, then
$(g')^a=(z^2)^a=1$, whence $a=0$.
This proves the isomorphism with $\widehat\Z$.
Multiplication by two is injective there, proving
uniqueness of the specified square root.
The displayed odd-prime corrections retain the
full profinite group, not only its pro-two quotient.

\emph{Which central extension occurs.}
In the two-class case, the explicit torsion quotient
$\chi:E\to\widehat\Z$ of D.60 has
$\chi(g)=1$ and vanishes on $C$.
Thus $\chi(g')=1$.
Applying it to $t^2=(g')^\kappa$ gives
$\kappa=2\chi(t)$, so $\kappa_2$ is even.
The preceding split case gives the involution $w$.
It cannot be $F$-conjugate to $v$:
otherwise $g'\in C_F(w)$ would lie in a conjugate
of $C\subseteq\overline{F'}$, contradicting its
primitive abelianization. By D.58 it therefore
belongs to the other actual section class.

In the one-class case, the same contradiction
excludes every outside involution centralizing $g'$,
because all such involutions are $F$-conjugate to $v$.
Thus $\kappa_2$ must be odd.
Equations (D.62m)--(D.62n) give precisely the root
and centralizer in (D.62f). Conversely either an
outside centralizing involution or an outside
square root is an outside centralizer element,
so these implications are equivalences.

\emph{Full basis and extension maps.}
For the two-class case, put $X'=g'$ and $Y'=vw$.
Both lie in $F$, and $[g',w]=1$ gives
\[
 [X',Y']=[g',v]=c,\qquad
 \alpha(X')=Y'X'(Y')^{-1},\qquad
 \alpha(Y')=(Y')^{-1}.
                                                        \tag{D.62o}
\]
For the one-class case put $X'=zv$ and $Y'=vz$.
Again both lie in $F$, and
\[
 [X',Y']=[z^2,v]=c,\qquad
 \alpha(X')=Y',\qquad \alpha(Y')=X'.
                                                        \tag{D.62p}
\]
In both cases the full boundary test of D.61 proves
that $X\mapsto X'$, $Y\mapsto Y'$ is an automorphism
of the full rank-two free profinite group.
The displayed intertwining identities extend it
to the asserted isomorphism with $v_0\mapsto v$;
its inverse on the kernel is the inverse free-group
automorphism, and on an element $fv^\epsilon$ it is
that inverse applied to $f$, followed by
$v_0^\epsilon$. This proves the exact extension
inverse and retains both parity labels.

Conversely, an isomorphism in this normalized
portion of D.61 sends $X$ in the two-class standard
model, or $XY$ in the one-class standard model,
to an element $g'$ with $[g',v]=c$.
The set of all such elements is precisely $gC$:
equality of the two conjugations of $v$ says that
$g^{-1}g'$ centralizes $v$.
The image of the standard other reflection in the
first model, or of its element $Xv_0$ in the second,
supplies the centralizer element in (D.62c).
This proves both directions of the final assertion.

The centralizer calculation and root formula give
no proof that $\alpha(g')$ is $F$-conjugate to $g'$
for any member of this coset. The continuation below
retains every unit multiplier and the other reflection
label through explicit additional transporter maps.
\end{proof}

\paragraph{Every unit multiplier and both actual reflection labels.}
The preceding criterion extends to the full parameter spaces
of D.61 by explicit transporter maps. No peripheral marking
is rescaled. Retain $c,v,g,C$ from (D.62b), and put
\[
 \begin{gathered}
 v_i=c^iv,\qquad \alpha_i=\operatorname{Inn}(v_i)|_F,
 \qquad C_i=C_F(v_i),\qquad i\in\{0,1\},\\
 C_0=C,\qquad C_1=gCg^{-1}.
 \end{gathered}                                         \tag{D.62q}
\]
For each actual unit $\mu\in\widehat\Z^\times$, define
\[
 g_{\mu,0}=c^{(\mu-1)/2}g,\qquad
 g_{\mu,1}=c^{(\mu+1)/2}g^{-1}.
                                                        \tag{D.62r}
\]
Both exponents are elements of $\widehat\Z$:
the dyadic coordinate of $\mu$ is odd, and division
by two is available at every odd prime.
The conjugation identities
\[
 c^avc^{-a}=c^{2a}v,\qquad
 gvg^{-1}=cv,\qquad g^{-1}(cv)g=v
\]
give, with no change of the named cusp,
\[
 \begin{gathered}
 g_{\mu,i}v_i g_{\mu,i}^{-1}=c^\mu v_i,\\
 \{a\in F:[a,v_i]=c^\mu\}=g_{\mu,i}C_i,\\
 [a]_{\rm ab}=(-1)^i[g]_{\rm ab}
           \quad(a\in g_{\mu,i}C_i).
 \end{gathered}                                         \tag{D.62s}
\]
The coset equality follows by dividing two transporters.
The homology equality follows from
$c\in\overline{F'}$ and
$C_i\subseteq\overline{F'}$, the latter being (D.59a)
and its displayed conjugate. In particular every one
of these transporters is homologically primitive;
its exact homology is independent of the unit $\mu$.

For a specified $a\in g_{\mu,i}C_i$ we now have
\[
 \alpha_i(a)=c^{-\mu}a,\qquad
 \alpha_i(a)\sim_F a
       \ \Longleftrightarrow\
                 C_E(a)\not\subset F.                  \tag{D.62t}
\]
The equivalence is the exact map $h\mapsto hv_i$
of (D.62j). In the two-class case, D.60 gives
$\chi(a)=(-1)^i$. Thus $t^2=a^\kappa$ for an
outside centralizer element implies
$2\chi(t)=(-1)^i\kappa$.
The even-dyadic argument (D.62k)--(D.62l) supplies
the unique outside commuting involution $w$.
It is in the other full section class because
$a$ has nonzero primitive homology whereas
every cusp-class fixed group lies in
$\overline{F'}$. The exact basis and extension maps are
\[
 X\longmapsto a,\qquad Y\longmapsto v_iw,\qquad
 v_0\longmapsto v_i,\qquad [X,Y]\longmapsto c^\mu.
                                                        \tag{D.62u}
\]
In the one-class case, that same fixed-group
argument excludes every outside involution in
$C_E(a)$. The odd-dyadic construction
(D.62m)--(D.62n) gives the unique $z\in E\setminus F$
with $z^2=a$, and the exact maps are
\[
 X\longmapsto zv_i,\qquad Y\longmapsto v_i z,\qquad
 v_0\longmapsto v_i,\qquad [X,Y]\longmapsto c^\mu.
                                                        \tag{D.62v}
\]
In both cases D.61's boundary test proves full
invertibility because $\mu$ is a full unit.
The inverse extension maps are the inverse kernel
map together with the specified label-$1$ reflection,
as in (D.62o)--(D.62p).

Conversely, every normalized conjugator in D.61
has some specified $\mu$ and $i$. In the two-class
model its image of $X$, and in the one-class model
its image of $XY$, belongs to the precise coset
in (D.62s). The image of the standard other
reflection or of $Xv_0$ supplies the outside
centralizer element, respectively. Therefore
\[
 \mathcal T_j(b)\ne\varnothing
 \quad\Longleftrightarrow\quad
 \left\{\begin{array}{l}
 \text{there are }\mu\in\widehat\Z^\times,\ i\in\{0,1\},
       \ a\in g_{\mu,i}C_i\\
 \text{with }c^{-\mu}a\sim_F a .
 \end{array}\right.
                                                        \tag{D.62w}
\]
Here $\mathcal T_j(b)$ is the exact peripheral
conjugator space defined in D.61, with its standard
model chosen by the section count. Equations
(D.62u)--(D.62v) and the converse give the actual
maps proving this equivalence. They retain the
unit exponent $c^\mu$ and both reflection labels
throughout. The existence in the right-hand set
remains unresolved; none of these identities
asserts it for an arbitrary full negative action.

The executable
\path{code/part_c/torus_unit_transporters_d62.py}
checks both reflection-label identities in the full
free-product word normal form for thirty-two odd
integer exponents. These identities hold before
imposing the boundary condition; the full unit
hypothesis is retained where D.61 proves invertibility.
The full centralizer and root statements are proved
above by their exact subgroup and coefficient maps.
